\section{Experimental design}
\label{sec:experiments}
\subsection{Development protocol and landmark evaluation}
The 801 development patients were partitioned into five patient-level folds. Fold 1 used 640 training and 161 validation patients; folds 2--5 each used 641 training and 160 validation patients. The same 175-patient holdout was evaluated by every fold model and every variant. The five fold-specific test evaluations are repeated evaluations of models trained on overlapping development partitions and tested on the same patients. They do not constitute five independent clinical cohorts. Patient identities were disjoint between training, validation, and holdout within the current split.

At each training epoch, a patient contributed one randomly selected available T0, T1, or T2 prefix. Only MRI observations available at or before the selected landmark were classifier inputs. Later observed visits, including T3, could supply latent targets for the auxiliary loss on training patients. Missing targets were masked. The 23 patients with only T0 remained eligible for pCR training but supplied no future-state loss. At evaluation, all three landmarks were evaluated for the 175 holdout patients, each of whom had T0, T1, and T2. T3 was not an evaluation landmark and was not used as observed input for these predictions.

\subsection{Model comparisons}
The principal comparison comprises three variants trained with seed 44 using the same five development folds. Every fold model is evaluated separately. Performance is summarized by the equal mean of the five fold-specific metric values. \nojepa{} uses direct pCR readout and a zero future-state loss weight. \direct{} uses the same direct readout with a future-state loss weight of 0.5. \fullmodel{} retains that auxiliary weight and replaces the direct readout with a concatenation of direct pCR features, attention-pooled predicted terminal state, and attention-pooled predicted change. All variants retain the shared image encoder, clinical inputs, assigned-arm encoding, fusion architecture, and training schedule.

The \direct{}--\nojepa{} contrast most directly tests adding future-state supervision. The \fullmodel{}--\direct{} contrast jointly changes the use of forecast features, spatial pooling, and classifier input width. It cannot isolate attention alone. For that purpose, the supplementary \meanmodel{} comparison uses forecast readout with mean spatial pooling and otherwise matches the forecast-attention configuration. All manuscript-facing pCR results first compute metrics separately for each fold model on the fixed holdout and then average those metrics across folds. Probabilities and thresholded predictions are never averaged between fold models. No VQ-latent reconstruction-head, change-weighted-loss, or state-conditioned-head variants are included.

\subsection{Implementation and optimization}
The network uses a model width of 768, eight layers in each of the spatial encoder, multimodal fusion module, and predictor, 16 attention heads with 64 dimensions per head, and feed-forward width 3,072. The 1,024-dimensional concatenated attention-head representation is projected back to the 768-dimensional residual stream. The spatial token grid is $2\times4\times4$. The configuration uses eight state queries, four pCR queries, dropout 0.3, and an EMA coefficient of 0.999. The pretrained convolutional encoder is fine-tuned jointly with the remaining network. Its training-patient provenance requires additional documentation, as discussed in Section~\ref{sec:limitations}.

We use AdamW \citep{adamw2019} with an initial base learning rate of $10^{-4}$ for both the encoder and the main network, weight decay 0.05, batch size 16, no gradient accumulation, and gradient-norm clipping at 1.0. Training uses BF16 mixed precision. A ten-epoch warm-up increases the learning rate from $10^{-5}$ to $10^{-4}$, followed by cosine decay to $10^{-6}$ at epoch 100. Training is capped at 200 epochs, with early stopping after 50 epochs without improvement. The best checkpoint is selected independently in each fold using validation patient-equal-prefix F1 at threshold 0.5. This is an F1 selection protocol, even though AUROC is a principal reporting measure. Deterministic execution was not enabled, so a fixed seed does not imply bitwise reproducibility across hardware or software changes.

\subsection{Metrics, aggregation, and patient-level uncertainty}
We report 14 complementary metrics: AUROC, non-interpolated average precision (AP), F1, accuracy, sensitivity, specificity, positive and negative predictive value (PPV/NPV), balanced accuracy, Matthews correlation coefficient (MCC), Youden index, negative log-likelihood (NLL), Brier score, and expected calibration error (ECE) with ten equal-width probability bins. Confusion-matrix counts are additionally retained. Thresholded primary metrics use 0.5. MCC and balanced accuracy help interpret performance under class imbalance~\citep{chicco2020mcc}; positive prevalence is $56/175=0.32$. No holdout threshold optimization or probability recalibration was performed.

Let $p_{if\ell}$ be the probability from the seed-44 model trained in fold $f$ for patient $i$ at landmark $\ell$. Within each variant, we define
\begin{equation}
\begin{aligned}
 M_{f\ell}&=M\bigl(\{y_i,p_{if\ell}\}_{i=1}^{175}\bigr),\\
 M_{f,\mathrm{macro}}&=\frac{1}{3}\sum_{\ell=0}^{2}M_{f\ell},\\
 \widehat M_\ell&=\frac{1}{5}\sum_{f=1}^{5}M_{f\ell},\\
 s_{M,\ell}&=\sqrt{\frac{1}{4}\sum_{f=1}^{5}
 (M_{f\ell}-\widehat M_\ell)^2}.
\end{aligned}
\end{equation}
Macro is first computed within each fold by averaging T0--T2 metric values, including F1, and is then averaged over folds. The corresponding Macro standard deviation is calculated from the five fold-specific Macro values. Main tables report the fold mean and sample standard deviation, with denominator four. This spread describes fitted-model differences associated with the five overlapping training partitions under seed 44. It is not a training-seed standard deviation, standard error, or patient-level interval. All probabilities remain specific to individual models; the reported mean is not the performance of a probability ensemble or of a model trained on the full development cohort.

For conditional patient-level uncertainty, we draw 2,000 bootstrap samples of the 175 patients with replacement. Every draw is shared across all fold models, variants, and landmarks. Within each draw, we recompute metrics separately for each fold model, average landmark metrics for Macro, and then equally average the five fold metrics. The 2.5th and 97.5th percentiles give unadjusted intervals for the fold-mean metric or paired difference. The fitted models are held fixed and folds are not resampled. These intervals exclude retraining, seed selection, model selection, and aggregation-policy uncertainty. Neither fold count nor repeated landmarks increase the independent patient count beyond 175.

ROC and precision--recall curves are computed separately for each fold model, interpolated onto common false-positive-rate and recall grids, respectively, and then averaged for display. The numerical AUROC and non-interpolated AP are computed per fold before averaging; their means need not equal the area under the displayed mean curve. Pointwise bands use shared patient resamples. Reliability diagrams form five equal-count rank groups of 35 patients separately within each fold model and then average group summaries across folds. These are descriptive summaries of several predictors, distinct from the ten equal-width bins used to calculate each fold model's ECE.

\subsection{Exploratory paired significance checks}
We use approximate, two-sided, null-centered paired patient-bootstrap tests for all 14 metrics, including AUROC~\citep{davison1997}. The tested quantity is the difference in the metric averaged across the five fixed fold models. We generate $B=99{,}999$ shared patient resamples with random-number seed 20261103. Every resample draws 175 patients with replacement and is shared across folds, models, and landmarks. Metrics are recomputed separately for each fold before equally averaging, preserving the repeated-patient dependence. No model is retrained, no fold is resampled, and no fold-level $t$ test is used.

Let $\Delta_b^*$ denote the fold-mean metric difference in resample $b$ and $\Delta_{\mathrm{obs}}$ its observed value. We estimate the two-sided $p$ value by
\begin{equation}
 p=\frac{1+\sum_{b=1}^{B}\mathbb{I}\bigl(|\Delta_b^*-\Delta_{\mathrm{obs}}|\geq|\Delta_{\mathrm{obs}}|\bigr)}{B+1}.
\end{equation}
Ties are included. The plus-one convention prevents Monte Carlo zero values, but does not make this bootstrap approximation finite-sample exact. Non-smooth statistics such as AP, F1, and binned ECE warrant particular caution.

The principal reporting family contains nine tests: the three pairwise comparisons among \nojepa{}, \direct{}, and \fullmodel{}, each for Macro AUROC, AP, and F1. The secondary family contains 81 tests: 27 landmark-specific AUROC/AP/F1 comparisons among the principal models; 33 comparisons of the other 11 Macro metrics; 12 \fullmodel{}--\meanmodel{} comparisons for AUROC/AP/F1 at T0, T1, T2, and Macro; and nine T2--T0 AUROC/AP/F1 comparisons within the three principal models. Holm adjustment~\citep{holm1979} is applied separately within each family. The aggregate CSV additionally reports an all-90 Holm sensitivity analysis; separate-family adjustments do not control the familywise error rate over the union. We use $p_{\mathrm{Holm}}<0.05$, with $*$, $**$, and $***$ denoting adjusted values below 0.05, 0.01, and 0.001, respectively, and ns otherwise.

The 2,000-resample percentile intervals remain unadjusted and are not exact inversions of these tests. We compute $p$ values from the bootstrap differences rather than deriving them from whether an interval includes zero. Seed 44, the metric-averaging policy, testing procedure, and reporting families were chosen after results were available. All inference is exploratory and conditions on these fitted models; it does not measure uncertainty from retraining or retrospective selection. Validation-threshold, subgroup, decision-curve, and training-probe summaries receive no additional significance tests.

\subsection{Operating-point and hypothetical decision analyses}
As a threshold sensitivity analysis, we select a separate Youden threshold for each variant, fold, and landmark using that fold model's validation predictions. The highest threshold resolves ties. Validation coverage per fold is 160--161 patients at T0, 151--154 at T1, and 138--143 at T2. We then apply the threshold to holdout predictions from that same fold model. We do not combine validation scores across folds, and no holdout outcome selects a threshold. This keeps the validation and test score definitions tied to one fitted model. Validation patients also informed that model's best-checkpoint selection, however, so this is not nested or independent calibration. The resulting holdout metrics are averaged across folds using the same procedure as the fixed-threshold metrics. Patient-bootstrap intervals keep the validation-selected thresholds fixed and exclude their selection uncertainty.

We additionally compute exploratory decision curves~\citep{vickers2006dca}. At risk threshold $q$, the generic positive action is applied when predicted pCR probability is at least $q$. Its mathematical net benefit is
\begin{equation}
 \mathrm{NB}(q)=\frac{\mathrm{TP}(q)}{n}
 -\frac{\mathrm{FP}(q)}{n}\frac{q}{1-q}.
\end{equation}
We calculate this separately from each fold model's probabilities over thresholds 0.05--0.80, average the three landmark net benefits for Macro within each fold, and then average across the five folds. Treat-all and treat-none are references. The calculation uses uncalibrated probabilities and does not validate a particular escalation, de-escalation, or other management action.

\subsection{Exploratory subgroup and mechanism analyses}
We examine released HR/HER2 combinations, age below or at least 50 years, native source namespace, and saved predicted-component crop coverage. Age 50 is an exploratory reporting cutpoint. Subgroup AUROC/AP are suppressed when fewer than ten patients or fewer than five outcomes in either class are available. Subgroup metrics are computed separately for each fold model and then equally averaged. Paired within-subgroup patient resampling preserves fold, model, and landmark dependence; single-class draws are omitted for AUROC intervals and valid-draw counts retained. Intervals are unadjusted, and no subgroup-interaction hypothesis is tested. Source namespaces share the same clinical release and are not independent external cohorts.

Mechanism analysis is restricted to archived fixed-training-probe diagnostics collected every five epochs. These probes sample up to 32 training patients and are not full validation or holdout forecast evaluations. Unlike the missing-target-masked training objective, the archived forecast diagnostic includes only continuous observed visit chains: the real landmark, target, and every intermediate canonical visit must be available. A gap excludes that landmark--target pair, and eligible-patient counts therefore vary by horizon. We report a descriptive common epoch 20 and common-epoch trajectories. Forecast skill is $1-\mathrm{L1}_{\mathrm{forecast}}/\mathrm{L1}_{\mathrm{copy}}$ within the same snapshot's EMA representation; the copy forecast holds the observed landmark state constant. This normalization avoids comparing raw latent errors between separately learned teacher coordinates. Predictor gradient norms from the pCR and weighted rollout objectives are descriptive checks of gradient routing, not independent statistical samples or clinical effects. These archived diagnostics are computed within each seed-44 fold model's own EMA coordinates and summarized descriptively across the five folds. Only epochs with complete diagnostic coverage are used for common-epoch summaries. Latent states are not combined between models. These training diagnostics are distinct from the fold-specific holdout pCR evaluation.

\subsection{Scope of the selected-seed fold-mean analysis}
Seed 44 and the decision to report a mean of fold-specific metrics were selected retrospectively after holdout results were available. Applying the same seed, folds, and averaging rule to every model preserves a matched architectural comparison, but does not remove selection bias. The fold standard deviation describes five fitted models trained on overlapping partitions under one seed; it does not characterize training-seed robustness. Patient-bootstrap inference conditions on those fixed models. A confirmatory study should lock training seeds, aggregation, model selection, statistical comparisons, and operating-point policies before evaluating independent patients.
